We set out to determine whether the coalescence phase and polarization angle of compact-binary signals are recoverable from strain by direct neural regression, individually or in their physically motivated sum/difference combinations.
The answer is \textbf{no}.
Every configuration tested converged to the optimal constant or noise-like predictor at the analytic random-guessing error on both angles, while the same shared representation recovered chirp mass, merger time, signal-to-noise ratio, and --- for one architecture --- a genuine partial sky-localization skill, from the same features.
No model here passed the pre-registered metric-health gate; the null instead rests on convergence across direct collapse inspection, a gradient-health trace, a formal permutation bootstrap, and population stratification, applied uniformly across every down-selected model rather than to a certified subset --- a defensible basis for the conclusion, argued explicitly in \S\ref{sec:phicpsi:reconcile}, but one that has not yet been through independent adversarial scrutiny.
Two findings exist only because this study looked further than strictly necessary: \cnnAttn's real, SNR-verified sky-localization skill, and inclination's small-but-real departure from the clean-noise-floor role this study's own methodology initially assumed for it.
Neither changes the central verdict; both sharpen, honestly, what it does and does not claim.
The degeneracy remains, on the evidence presented here, effectively exact for this population as a point-estimation problem.
